\title{: Advancing the Boundaries of Mathematical Reasoning in Open Language Models}

\begin{abstract}
Mathematical reasoning presents a notable challenge for language models due to its intricate and structured characteristics. In this work, we present DeepSeekMath~7B, which extends the pre-training of DeepSeek-Coder-Base-v1.5 7B by incorporating 120B math-related tokens drawn from Common Crawl, alongside natural language and code datasets. DeepSeekMath~7B attains a remarkable 51.7\% score on the competition-level MATH benchmark without utilizing external tools or ensemble voting methods, nearing the performance of Gemini-Ultra and GPT-4. Applying self-consistency with 64 samples from DeepSeekMath~7B yields 60.9\% on MATH. The strong mathematical reasoning skills of DeepSeekMath~are credited to two main contributions: first, leveraging the vast potential of publicly accessible web data through a carefully designed data selection pipeline; second, proposing Group Relative Policy Optimization (GRPO), a PPO variant that improves mathematical reasoning capabilities while simultaneously reducing PPO’s memory consumption.
\end{abstract}

\section{Introduction}

Large language models (LLMs) have transformed the methodology for tackling mathematical reasoning within AI, driving substantial progress on quantitative reasoning benchmarks \citep{MATH} as well as geometry reasoning benchmarks \citep{trinh2024solving}. Furthermore, these models have demonstrated their value in assisting humans with solving challenging mathematical problems \citep{tao}. Nonetheless, state-of-the-art models like GPT-4 \citep{gpt4} and Gemini-Ultra \citep{gemini} remain inaccessible to the public, and the currently available open-source models lag considerably behind in terms of performance.

In this work, we present DeepSeekMath, a language model tailored specifically for the mathematics domain that markedly surpasses the mathematical proficiency of existing open-source models and nears GPT-4’s performance on academic evaluations. To accomplish this, we compile the DeepSeekMath Corpus, a substantial, high-quality pre-training dataset containing 120B math-related tokens. This corpus is derived from Common Crawl (CC) data using a classifier based on fastText \citep{joulin2016fasttext}. Initially, the classifier is trained with positive samples drawn from OpenWebMath \citep{openwebmath}, while negative samples are gathered from a broad variety of other web pages. Following this, we utilize the classifier to identify additional positive samples within CC, which are then refined through human annotation. The classifier is subsequently retrained with this enhanced dataset to boost its accuracy. Evaluation results demonstrate the high quality of the large-scale corpus, as our base DeepSeekMath-Base 7B model attains 64.2\% accuracy on GSM8K \citep{gsm8k} and 36.2\% on the advanced MATH dataset \citep{MATH}, surpassing the performance of Minerva 540B \citep{minerva}. Moreover, since the DeepSeekMath Corpus is multilingual, we observe improvements on Chinese mathematics benchmarks \citep{wei2023cmath,agieval}. We view our methodology for processing mathematical data as an initial framework for the research community, with ample opportunities for future enhancements.

DeepSeekMath-Base is initialized using DeepSeek-Coder-Base-v1.5 7B \citep{deepseek-coder}, as we find that beginning with a model trained on code is preferable to starting from a general large language model. Additionally, we observe that math training not only enhances the model's performance on the MMLU \citep{mmlu} and BBH benchmarks \citep{bbh}, but also strengthens its overall reasoning abilities beyond mathematics.

Following pre-training, we conduct mathematical instruction tuning on DeepSeekMath-Base utilizing data involving chain-of-thought \citep{cot}, program-of-thought \citep{pot,pal}, and tool-integrated reasoning \citep{tora}. The resulting model, DeepSeekMath-Instruct 7B, outperforms all other 7B-scale models and achieves performance comparable to 70B open-source instruction-tuned models.

Moreover, we propose Group Relative Policy Optimization (GRPO), a variant of the reinforcement learning (RL) algorithm Proximal Policy Optimization (PPO) \citep{schulman2017proximal}. GRPO eliminates the need for a critic model by estimating the baseline from group-level scores, substantially reducing the computational resources required for training. Using only a portion of English instruction tuning data, GRPO achieves significant improvements over the already strong DeepSeekMath-Instruct, both on in-domain mathematical benchmarks (GSM8K: 82.9\% $\rightarrow$ 88.2\%, MATH: 46.8\% $\rightarrow$ 51.7\%) and out-of-domain math tasks (e.g., CMATH: 84.6\% $\rightarrow$ 88.8\%) during the reinforcement learning process. We also present a unified framework to interpret various approaches, such as Rejection Sampling Fine-Tuning (RFT) \citep{yuan2023scaling}, Direct Preference Optimization (DPO) \citep{dpo}, PPO, and GRPO. Within this unified framework, we show that these methods can be viewed as either direct or simplified reinforcement learning techniques. We further perform comprehensive experiments—comparing online versus offline training, outcome versus process supervision, single-turn versus iterative RL, and more—to thoroughly examine the fundamental components of this framework. Finally, we discuss why our reinforcement learning approach enhances the performance of instruction-tuned models and outline future directions for achieving more effective RL based on this unified perspective.

\subsection{Contributions}
Our work presents scalable mathematical pre-training alongside the investigation and evaluation of reinforcement learning techniques.

\textbf{Math Pre-Training at Scale}

\begin{itemize}[topsep=0pt]
    \item We demonstrate strong evidence that publicly available Common Crawl data encompasses valuable mathematical content. Utilizing a carefully crafted data selection pipeline, we build the DeepSeekMath Corpus, a high-quality dataset containing 120B tokens sourced from web pages curated for mathematical material. This corpus is nearly 7 times larger than the math web pages used by Minerva \citep{minerva} and 9 times larger than the recently introduced OpenWebMath \citep{openwebmath}.

    \item Our pre-trained base model, DeepSeekMath-Base 7B, attains performance on par with Minerva 540B \citep{minerva}, suggesting that the number of parameters alone does not fully determine mathematical reasoning proficiency. A more compact model trained on high-quality data can also achieve strong results.

    \item We report insights from math training experiments. Pre-training on code before math training enhances the model's capacity to solve mathematical problems both with and without the use of external tools. This provides a partial resolution to the enduring question: \textit{does code training enhance reasoning abilities?} Our findings suggest that it does, at least in the context of mathematical reasoning.

    \item Despite the prevalent use of arXiv papers in many math-focused studies, training on this data does not yield significant improvements across all mathematical benchmarks considered in this work.
\end{itemize}

\textbf{Exploration and Analysis of Reinforcement Learning}

\begin{itemize}[topsep=0pt]
    \item We propose Group Relative Policy Optimization (GRPO), a reinforcement learning algorithm that is both efficient and effective. GRPO eliminates the need for a critic model by estimating the baseline from group scores, which greatly reduces the computational resources required compared to Proximal Policy Optimization (PPO).
    \item We show that GRPO substantially improves the performance of our instruction-tuned model DeepSeekMath-Instruct using only instruction-tuning data. Additionally, we note improvements in out-of-domain performance during the reinforcement learning phase.
    \item We present a unified framework to interpret various approaches, including RFT, DPO, PPO, and GRPO. We also perform thorough experiments—such as comparing online versus offline training, outcome versus process supervision, and single-turn versus iterative reinforcement learning—to deeply analyze the core components of this framework.
    \item Building on this unified framework, we investigate the underlying reasons behind reinforcement learning’s effectiveness and outline several promising directions to enhance reinforcement learning for LLMs.
\end{itemize}

\subsection{Summary of Evaluations and Metrics}

\begin{itemize}[topsep=0pt]
    \item \textbf{English and Chinese Mathematical Reasoning}: We carry out extensive evaluations of our models on English and Chinese benchmark datasets spanning mathematical problems from elementary to college level. English benchmarks comprise GSM8K \citep{gsm8k}, MATH \citep{MATH}, SAT \citep{llemma}, OCW Courses \citep{minerva}, and MMLU-STEM \citep{mmlu}. Chinese benchmarks include MGSM-zh \citep{mgsm}, CMATH \citep{wei2023cmath}, Gaokao-MathCloze \citep{agieval}, and Gaokao-MathQA \citep{agieval}. We assess models’ capabilities to generate self-contained textual solutions without tool assistance, as well as their capacity to solve problems utilizing Python.
    
    On English benchmarks, DeepSeekMath-Base competes closely with the closed-source Minerva 540B \citep{minerva}, outperforming all open-source base models (e.g., Mistral 7B \citep{mistral} and Llemma-34B \citep{llemma}), irrespective of whether those models have undergone math pre-training, often by a notable margin. Importantly, DeepSeekMath-Base achieves superior results on Chinese benchmarks, likely because, unlike prior works \citep{minerva,llemma}, we incorporate high-quality non-English mathematical pre-training data rather than relying solely on English data. With mathematical instruction tuning and reinforcement learning, our DeepSeekMath-Instruct and DeepSeekMath-RL models exhibit strong capabilities, achieving an accuracy exceeding 50\% on the challenging MATH dataset—a first in the open-source domain.
\end{itemize}

\begin{itemize}
    \item \textbf{Formal Mathematics}: We assess DeepSeekMath-Base by applying it to the informal-to-formal theorem proving task introduced in \citep{dsp_proof} on miniF2F \citep{minif2f}, using Isabelle \citep{isabelle} as the proof assistant. DeepSeekMath-Base shows strong few-shot autoformalization capabilities.
    \item \textbf{Natural Language Understanding, Reasoning, and Code}: To establish a thorough evaluation of the models' general comprehension, reasoning, and coding skills, we test DeepSeekMath-Base on the Massive Multitask Language Understanding (MMLU) benchmark \citep{mmlu}, which includes 57 multiple-choice tasks across various subjects, on BIG-Bench Hard (BBH) \citep{bbh}, comprising 23 demanding tasks primarily requiring multi-step reasoning, and on HumanEval \citep{codex} and MBPP \citep{mbpp}, which are standard benchmarks for assessing code language models. Math pre-training enhances both language understanding and reasoning capabilities.
\end{itemize}

\section{Math Pre-Training}

\subsection{Data Collection and Decontamination} In this section, we describe the method used to build the DeepSeekMath Corpus from Common Crawl. As illustrated in Figure \ref{fig:data_collect}, we outline an iterative pipeline detailing how to systematically extract a large-scale mathematical dataset from Common Crawl, starting with a seed corpus (for example, a small but high-quality math-related dataset). It is important to note that this methodology can be applied to other fields as well, such as programming.

Initially, we select OpenWebMath \citep{openwebmath}, a collection of high-quality mathematical web documents, as our seed corpus. Using this dataset, we train a fastText model \citep{joulin2016fasttext} to retrieve web pages similar to those in OpenWebMath. Specifically, we randomly sample 500,000 entries from the seed corpus as positive training samples and 500,000 web pages from Common Crawl as negative samples. The training is performed using an open-source library\footnote{\url{https://fasttext.cc}}, with the vector dimension set to 256, learning rate at 0.1, maximum word n-gram length of 3, minimum word occurrence count of 3, and number of training epochs fixed at 3. To reduce the volume of the original Common Crawl data, we apply URL-based deduplication and near-duplicate removal, resulting in 40 billion HTML web pages. We then use the fastText model to recall mathematical web pages from this deduplicated Common Crawl dataset. To eliminate low-quality mathematical content, we rank the retrieved pages based on their fastText scores and retain only the highest-ranked pages. The amount of data retained is determined through pre-training experiments conducted on the top 40B, 80B, 120B, and 160B tokens. In the initial iteration, we opt to keep the top 40B tokens.

Following the initial round of data collection, a significant number of mathematical web pages remain uncollected, primarily because the fastText model was trained on a set of positive examples that lacked sufficient diversity. To address this, we identify additional mathematical web sources to expand the seed corpus, aiming to enhance the fastText model. Concretely, we begin by partitioning the entire Common Crawl dataset into distinct domains; each domain is defined as a group of web pages sharing the same base URL. For each domain, we compute the proportion of web pages collected during the first iteration. Domains in which more than 10\% of the web pages were gathered are categorized as math-related (e.g., \url{mathoverflow.net}).

Next, we manually label the URLs connected to mathematical content within these identified domains (for example, \url{mathoverflow.net/questions}). Web pages associated with these URLs that were not yet collected are added to the seed corpus. This method allows us to acquire additional positive examples, thereby training a more effective fastText model that can retrieve a greater volume of mathematical content in subsequent iterations. After four rounds of data collection, the dataset comprises 35.5M mathematical web pages, totaling 120 billion tokens. Notably, in the fourth iteration, we observe that almost 98\% of the data had already been captured by the third iteration, prompting us to halt further data collection.

To prevent benchmark contamination, we adopt the filtering strategy from \cite{deepseek-coder} to exclude web pages containing questions or answers from English mathematical benchmarks such as GSM8K \citep{gsm8k} and MATH \citep{MATH}, as well as Chinese benchmarks including CMATH \citep{wei2023cmath} and AGIEval \citep{agieval}. Our filtering rule is as follows: any text segment containing a 10-gram string that exactly matches any substring from the evaluation benchmarks is removed from our math training corpus. For benchmark texts shorter than 10 grams but containing at least 3 grams, we perform exact matching to filter out contaminated web pages.

\subsection{Validating the Quality of the DeepSeekMath~Corpus}

We conduct pre-training experiments to evaluate how the DeepSeekMath~Corpus compares to recently released mathematics training corpora:

\begin{itemize}[topsep=0pt]
    \item \textbf{MathPile} \citep{mathpile}: a multi-source corpus comprising 8.9B tokens, aggregated from textbooks, Wikipedia, ProofWiki, CommonCrawl, StackExchange, and arXiv, with the majority (over 85\%) derived from arXiv;
    \item \textbf{OpenWebMath} \citep{openwebmath}: CommonCrawl data filtered specifically for mathematical content, amounting to 13.6B tokens;
    \item \textbf{Proof-Pile-2} \citep{llemma}: a mathematical corpus consisting of OpenWebMath, AlgebraicStack (which contains 10.3B tokens of mathematical code), and arXiv papers (28.0B tokens). In experiments involving Proof-Pile-2, we follow the approach of \cite{llemma} and use an arXiv:Web:Code token ratio of 2:4:1.
\end{itemize}

\subsubsection{Training Configuration}
\label{sec:quality-policy}
We perform math training on a general pre-trained language model consisting of 1.3B parameters, which follows the same architecture as the DeepSeek LLMs \citep{deepseek-llm}, referred to as DeepSeek-LLM 1.3B. Separate training is conducted for each mathematical corpus using 150B tokens. All experiments utilize the efficient and lightweight HAI-LLM \citep{haillm} training framework. Consistent with the training protocol of DeepSeek LLMs, we employ the AdamW optimizer \citep{adamW} with parameters $\beta_1=0.9$, $\beta_2=0.95$, and $\mathrm{weight\_decay}=0.1$, alongside a multi-phase learning rate schedule where the learning rate peaks after 2,000 warmup steps, then declines to 31.6\% after 80\% of training, and further reduces to 10.0\% of the peak following 90\% of the training process. The maximum learning rate is set to 5.3e-4, and the batch size is fixed at 4M tokens with a context length of 4K.

\subsubsection{Evaluation Outcomes}

\textbf{The DeepSeekMath~Corpus exhibits high quality, supports multilingual mathematical content, and is the largest in scale.}

\begin{itemize}[topsep=0pt]
    \item \textbf{High-quality}:
    We assess downstream performance on 8 mathematical benchmarks employing few-shot chain-of-thought prompting \cite{cot}. As presented in Table \ref{tab:corpora_comparison}, the model trained on the DeepSeekMath~Corpus clearly outperforms others. Figure \ref{fig:corpora_comparisons} demonstrates that at 50B tokens (equivalent to one full epoch of Proof-Pile-2), the model trained on DeepSeekMath consistently surpasses the performance of Proof-Pile-2, suggesting that the average quality of the DeepSeekMath Corpus is superior.
    \item \textbf{Multilingual}:
    The DeepSeekMath~Corpus includes data in various languages, with English and Chinese being the most prominent. Table \ref{tab:corpora_comparison} indicates that training on the DeepSeekMath~Corpus improves mathematical reasoning capabilities in both English and Chinese. In contrast, other existing mathematical corpora, which are predominantly English-focused, show limited gains and may even negatively impact Chinese mathematical reasoning performance.
    \item \textbf{Large-scale}:
    The DeepSeekMath~Corpus surpasses existing mathematical corpora by several folds in size. Figure \ref{fig:corpora_comparisons} illustrates that DeepSeek-LLM 1.3B trained on DeepSeekMath~Corpus follows a steeper learning trajectory and achieves more sustained improvements. Conversely, the baseline corpora are significantly smaller, having undergone multiple repetitions during training, with model performance quickly reaching a plateau.
\end{itemize}

\subsection{Training and Evaluation of DeepSeekMath-Base 7B}

This section presents DeepSeekMath-Base 7B, a foundational model designed with enhanced reasoning capabilities, particularly in mathematics. The model is initially based on DeepSeek-Coder-Base-v1.5 7B \citep{deepseek-coder} and trained over 500 billion tokens. The data composition consists of 56\% from the DeepSeekMath Corpus, 4\% from AlgebraicStack, 10\% from arXiv, 20\% sourced from GitHub code, and the remaining 10\% derived from natural language data in Common Crawl in both English and Chinese. The training configuration largely follows the setup described in Section \ref{sec:quality-policy}, except that the peak learning rate is adjusted to 4.2e-4 and the batch size is set to 10 million tokens.

We perform an extensive evaluation of the mathematical reasoning abilities of DeepSeekMath-Base 7B, emphasizing its capacity to generate complete mathematical solutions independently without external aids, tackle mathematical problems using tools, and perform formal theorem proving. Additionally, we assess the model's broader capabilities in natural language understanding, logical reasoning, and programming proficiency.

\paragraph{Mathematical Problem Solving via Stepwise Reasoning}

We assess DeepSeekMath-Base’s ability to solve mathematical questions using few-shot chain-of-thought prompting \citep{cot} across eight benchmarks in both English and Chinese. These benchmarks include quantitative reasoning datasets (such as GSM8K \citep{gsm8k}, MATH \citep{MATH}, and CMATH \citep{wei2023cmath}) and multiple-choice problems (like MMLU-STEM \citep{mmlu} and Gaokao-MathQA \citep{agieval}), spanning a wide range of mathematical topics from basic to university-level difficulty.

As demonstrated in Table \ref{tab:base_math_cot}, DeepSeekMath-Base 7B achieves the highest performance across all eight benchmarks among open-source base models, which includes the widely-adopted general model Mistral 7B \citep{mistral} and the newly introduced Llemma 34B \citep{llemma} that was trained in mathematics on Proof-Pile-2 \citep{llemma}. Remarkably, on the competition-level MATH dataset, DeepSeekMath-Base outperforms all existing open-source base models by an absolute margin exceeding 10\%, and even surpasses Minerva 540B \citep{minerva}, a closed-source base model that is 77 times larger, built upon PaLM \citep{palm} and further trained on mathematical literature.

\paragraph{Mathematical Problem Solving with Tool Use} We assess program-assisted mathematical reasoning on GSM8K and MATH using few-shot program-of-thought prompting \citep{pot,pal}. Models are prompted to solve problems by generating Python code, leveraging libraries such as \textit{math} and \textit{sympy} for complex calculations. The output of the executed program is then used as the final answer. As illustrated in Table \ref{tab:base_math_pot_proof}, DeepSeekMath-Base 7B exceeds the previous state-of-the-art Llemma 34B.

\paragraph{Formal Mathematics}

Automating formal proof generation is crucial for guaranteeing the correctness and dependability of mathematical proofs while improving efficiency, a focus that has gained momentum recently. We test DeepSeekMath-Base 7B on the informal-to-formal proving task from \citep{dsp_proof}, which involves generating a formal proof given an informal statement, its formal counterpart, and an informal proof. The evaluation is performed on miniF2F \citep{minif2f}, a benchmark for formal Olympiad-level mathematics, where a formal proof is produced in Isabelle for each problem using few-shot prompting. Following the approach in \cite{dsp_proof}, models generate proof sketches, which are then completed by the off-the-shelf automated prover Sledgehammer \citep{sledgehammer}. As indicated in Table \ref{tab:base_math_pot_proof}, DeepSeekMath-Base 7B shows strong capabilities in proof autoformalization.

\paragraph{Natural Language Understanding, Reasoning, and Code}

We measure the model's natural language understanding on MMLU \citep{mmlu}, reasoning on BBH \citep{bbh}, and coding ability on HumanEval \citep{codex} and MBPP \citep{mbpp}. As shown in Table \ref{tab:understanding_reasoning_code}, DeepSeekMath-Base 7B achieves substantial improvements over its predecessor, DeepSeek-Coder-Base-v1.5 \citep{deepseek-coder}, in MMLU and BBH, demonstrating the beneficial effects of math training on language comprehension and reasoning skills. Furthermore, by incorporating code tokens during continual training, DeepSeekMath-Base 7B preserves the coding performance of DeepSeek-Coder-Base-v1.5 on both coding benchmarks. Overall, DeepSeekMath-Base 7B markedly outperforms the general model Mistral 7B \citep{mistral} across the three reasoning and coding benchmarks.

\section{Supervised Fine-Tuning}

\subsection{SFT Data Preparation}

We create a mathematical instruction-tuning dataset that includes English and Chinese problems from various mathematical domains and different difficulty levels: each problem is paired with solutions presented in chain-of-thought (CoT) \citep{cot}, program-of-thought (PoT) \citep{pot,pal}, and tool-assisted reasoning formats \citep{tora}. The dataset contains a total of 776K training instances.

\begin{itemize}[topsep=0pt]
    \item \textbf{English mathematical datasets}: We annotate GSM8K and MATH problems with tool-assisted solutions and incorporate a subset of MathInstruct \citep{MathInstruct} along with the training portion of Lila-OOD \citep{lila}, where problems are addressed using CoT or PoT methods. Our English dataset spans multiple mathematical disciplines such as algebra, probability, number theory, calculus, and geometry.
    \item \textbf{Chinese mathematical datasets}: We gather Chinese K-12 math problems covering 76 subtopics, including linear equations, with solutions provided in both CoT and tool-integrated reasoning formats.
\end{itemize}

\subsection{Training and Evaluation of DeepSeekMath-Instruct 7B}

Here, we present DeepSeekMath-Instruct 7B, which is fine-tuned on mathematical instructions based on the DeepSeekMath-Base model. Training samples are randomly concatenated until reaching a maximum sequence length of 4K tokens. The model is trained for 500 steps with a batch size of 256 and a fixed learning rate of 5e-5.

We assess the mathematical capabilities of the models both without and with tool usage across four quantitative reasoning benchmarks in English and Chinese. Our model's performance is compared against state-of-the-art models available at the time:

\begin{itemize}[topsep=0pt]
    \item \textbf{Closed-source models} include: (1) the GPT family, where GPT-4 \citep{gpt4} and GPT-4 Code Interpreter~\footnote{\url{https://openai.com/blog/chatgpt-plugins\#\#code-interpreter}} are the most advanced, (2) Gemini Ultra and Pro \citep{gemini}, (3) Inflection-2 \citep{inflection-2}, (4) Grok-1~\footnote{\url{https://x.ai/model-card}}, as well as recent models launched by Chinese firms such as (5) Baichuan-3~\footnote{\url{https://www.baichuan-ai.com}}, and (6) the newest GLM-4~\footnote{\url{https://open.bigmodel.cn/dev/api\#glm-4}} from the GLM series \citep{glm}.
\end{itemize}

These models serve general purposes, with most having undergone multiple alignment processes.  
\item \textbf{Open-source models} comprise: general models such as (1) DeepSeek-LLM-Chat 67B \citep{deepseek-llm}, (2) Qwen 72B \citep{qwen}, (3) SeaLLM-v2 7B \citep{seallm}, and (4) ChatGLM3 6B \citep{chatglm3}, alongside models enhanced for mathematics, including (5) InternLM2-Math 20B~\footnote{\url{https://github.com/InternLM/InternLM-Math}}, which is built on InternLM2 and trained on math tasks followed by instruction tuning, (6) Math-Shepherd-Mistral 7B that employs PPO training \citep{schulman2017proximal} on Mistral 7B \citep{mistral} with a reward model guided by process supervision, (7) the WizardMath series \citep{wizardmath} that improves mathematical reasoning in Mistral 7B and Llama-2 70B \citep{llama2} through evolve-instruct (a variant of instruction tuning using AI-evolved instructions) and PPO training with primary training problems from GSM8K and MATH, (8) MetaMath 70B \citep{metamath}, which fine-tunes Llama-2 70B on an augmented GSM8K and MATH dataset, (9) ToRA 34B \cite{tora}, a CodeLlama 34B fine-tuned for tool-integrated mathematical reasoning, and (10) MAmmoTH 70B \citep{MathInstruct}, which instruction-tunes Llama-2 70B on MathInstruct.
\end{itemize}

As presented in Table \ref{tab:sft_rl_math}, in the evaluation setting where tool use is not permitted, DeepSeekMath-Instruct 7B exhibits strong step-by-step reasoning capabilities. Notably, on the competitive MATH dataset, our model outperforms all open-source models and most proprietary ones (e.g., Inflection-2 and Gemini Pro) by at least 9\% in absolute terms. This advantage holds even against substantially larger models (such as Qwen 72B) or those specifically enhanced via math-focused reinforcement learning (e.g., WizardMath-v1.1 7B). While DeepSeekMath-Instruct is competitive with Chinese proprietary models GLM-4 and Baichuan-3 on MATH, it still trails behind GPT-4 and Gemini Ultra.

In the evaluation setting that allows the use of natural language reasoning combined with program-based tool utilization for problem solving, DeepSeekMath-Instruct 7B reaches nearly 60\% accuracy on MATH, exceeding all existing open-source models. On other benchmarks, our model remains competitive with DeepSeek-LLM-Chat 67B, the previous state-of-the-art which is ten times larger.

\section{Reinforcement Learning}

\subsection{Group Relative Policy Optimization} Reinforcement learning (RL) has proven effective for enhancing the mathematical reasoning abilities of LLMs after the Supervised Fine-Tuning (SFT) phase \citep{wang2023math,wizardmath}. Here, we introduce our efficient and effective RL algorithm named Group Relative Policy Optimization (GRPO).

\subsubsection{From PPO to GRPO} Proximal Policy Optimization (PPO) \citep{schulman2017proximal} is a widely adopted actor-critic RL algorithm used during the RL fine-tuning of LLMs \citep{ouyang2022training}. Specifically, it improves LLMs by maximizing the following surrogate objective:

\begin{equation}
\footnotesize
    \mathcal{J}_{PPO}(\theta) = \mathbb{E}{[q \sim P(Q), o \sim \pi_{\theta_{old}}(O|q)]} \frac{1}{|o|} \sum_{t=1}^{|o|} \min \left[ \frac{\pi_\theta(o_{t} | q, o_{<t})}{\pi_{\theta_{old}}(o_{t} | q, o_{<t})} A_{t}, \text{clip} \left( \frac{\pi_\theta(o_{t} | q, o_{<t})}{\pi_{\theta_{old}}(o_{t} | q, o_{<t})}, 1 - \varepsilon, 1 + \varepsilon \right)  A_{t} \right] ,
\end{equation}

Here, $\pi_{\theta}$ and $\pi_{\theta_{old}}$ denote the current and previous policy models, respectively, while $q$ and $o$ represent questions and outputs sampled from the question dataset and the old policy $\pi_{\theta_{old}}$. The parameter $\varepsilon$ is a clipping-related hyper-parameter introduced in PPO to stabilize the training process. The advantage $A_t$ is computed using Generalized Advantage Estimation (GAE) \citep{gae}, relying on the rewards $\{r_{\ge t}\}$ and a learned value function $V_{\psi}$. Consequently, PPO requires training a value function alongside the policy model. To prevent excessive optimization of the reward model, the standard technique involves incorporating a per-token KL penalty from a reference model into the reward at each token \citep{ouyang2022training}, given by

\begin{equation}
    r_{t} = r_\varphi(q, o_{\le t}) - \beta \log\frac{\pi_{\theta}(o_{t}|q, o_{<t})}{\pi_{ref}(o_{t}|q, o_{<t})},
\label{eq:PPO-reward}
\end{equation}

where $r_\varphi$ is the reward model, $\pi_{ref}$ refers to the reference model, typically the initial SFT model, and $\beta$ is the coefficient controlling the strength of the KL penalty.

Since the value function in PPO is usually a model of similar size to the policy model, this adds significant memory and computational overhead. Moreover, during RL training, the value function is used as a baseline to reduce variance when computing the advantage. In the context of LLMs, the reward model usually assigns a reward score only to the final token, complicating the training of an accurate value function at every token. To tackle this, as illustrated in Figure \ref{fig:grpo}, we introduce Group Relative Policy Optimization (GRPO), which eliminates the need for an additional value function approximation as required in PPO. Instead, GRPO employs the average reward of multiple sampled outputs generated in response to the same question as the baseline. Specifically, for each question $q$, GRPO samples a set of outputs $\{o_1, o_2, \ldots, o_G\}$ from the old policy $\pi_{\theta_{old}}$, then optimizes the policy model by maximizing the objective:

\begin{equation}
\footnotesize
\begin{split}
    \mathcal{J}_{GRPO}(\theta) &= \mathbb{E}_{q \sim P(Q), \{o_i\}_{i=1}^G \sim \pi_{\theta_{old}}(O|q)}  \\
    & \frac{1}{G}\sum_{i=1}^G\frac{1}{|o_i|} \sum_{t=1}^{|o_i|} \left\{ \min \left[ \frac{\pi_\theta(o_{i,t} | q, o_{i,<t})}{\pi_{\theta_{old}}(o_{i,t} | q, o_{i,<t})} \hat{A}_{i,t}, \text{clip} \left( \frac{\pi_\theta(o_{i,t} | q, o_{i,<t})}{\pi_{\theta_{old}}(o_{i,t} | q, o_{i,<t})}, 1 - \varepsilon, 1 + \varepsilon \right)  \hat{A}_{i,t} \right] - \beta \mathbb{D}_{KL}\left[\pi_{\theta} || \pi_{ref}\right]\right\} ,
\end{split}
\label{eq:GRPO-obj}
\end{equation}

where $\varepsilon$ and $\beta$ are hyper-parameters, and $\hat{A}_{i,t}$ denotes the advantage computed solely from the relative rewards of outputs within each group, details of which will be provided in subsequent subsections. The group-relative strategy used by GRPO to calculate advantages aligns naturally with the comparative character of reward models, which are commonly trained on datasets containing comparisons among outputs for the same question. Additionally, unlike PPO, where the KL penalty is incorporated into the reward, GRPO applies regularization by directly adding the KL divergence between the trained policy and the reference policy into the loss function, simplifying the calculation of $\hat{A}_{i,t}$. Distinct from the KL penalty term in (\ref{eq:PPO-reward}), we estimate the KL divergence using the following unbiased estimator \citep{kl_approx}:

\begin{equation}
\small
    \mathbb{D}_{KL}\left[\

\begin{algorithm}[t]
  \small
  \caption{Iterative Group Relative Policy Optimization}
  \textbf{Input} initial policy model $\pi_{\theta_{\text{init}}}$; reward models $r_\varphi$; task prompts $\mathcal{D}$; 
  hyperparameters $\varepsilon$, $\beta$, $\mu$
  \begin{algorithmic}[1]
    \State Initialize policy model $\pi_\theta$ with $\pi_{\theta_{\text{init}}}$
    \For{iteration = 1, \dots, I}
       \State Set reference model $\pi_{ref}$ to current policy $\pi_{\theta}$
      \For{step = 1, \dots, M}
      \State Draw a batch $\mathcal{D}_b$ from dataset $\mathcal{D}$
      \State Save current policy as old policy: $\pi_{\theta_{old}} \leftarrow \pi_{\theta}$ 
      \State For each question $q \in \mathcal{D}_b$, sample $G$ outputs $\{o_i\}_{i=1}^G$ from $\pi_{\theta_{old}}(\cdot \mid q)$
      \State Evaluate the reward values $\{r_i\}_{i=1}^G$ for each output $o_i$ using $r_{\varphi}$
      \State Calculate the group relative advantage $\hat{A}_{i,t}$ for the $t$-th token in output $o_i$
      \For{GRPO iteration = 1, \dots, $\mu$}
        \State Update policy $\pi_{\theta}$ by maximizing the GRPO objective (refer to Equation \ref{eq:GC-GRPO})
      \EndFor
    \EndFor 
    \State Continuously update $r_\varphi$ using a replay buffer for training 
    \EndFor 
  \end{algorithmic}
  \textbf{Output} $\pi_\theta$
  \label{alg:iter-grpo}
\end{algorithm}

\subsubsection{Outcome Supervision RL with GRPO} In formal terms, for a given question $q$, a set of outputs $\{o_1, o_2, \cdots, o_G\}$ is drawn from the previous policy $\pi_{\theta_{old}}$. The reward model then assigns a score to each output, producing a reward vector $\mathbf{r}=\{r_1, r_2, \cdots, r_G\}$. These rewards are normalized by subtracting their mean and dividing by the group’s standard deviation. Outcome supervision provides the normalized reward at the conclusion of each output $o_i$ and sets the advantage values $\hat{A}_{i, t}$ for every token in $o_i$ to this normalized reward, i.e., $\hat{A}_{i, t} = \widetilde{r}_i = \frac{r_i - {\rm mean}(\mathbf{r})}{{\rm std}(\mathbf{r})}$. The policy is then optimized by maximizing the objective described in equation (\ref{eq:GRPO-obj}).

\subsubsection{Process Supervision RL with GRPO} Outcome supervision delivers a reward only at the end of each output, which might be inadequate and inefficient for guiding the policy in complex mathematical tasks. Following the approach in \cite{wang2023math}, we also investigate process supervision, which assigns rewards at the conclusion of each reasoning step. Concretely, given a question $q$ and $G$ sampled outputs $\{o_1, o_2, \cdots, o_G\}$, the process reward model evaluates each step within the outputs, generating rewards: $\mathbf{R} = \{ \{r_1^{index(1)},\cdots,r_1^{index(K_1)}\}, \cdots,  \{r_G^{index(1)},\cdots,r_G^{index(K_G)}\} \}$, where $index(j)$ indicates the token index marking the end of the $j$-th step, and $K_i$ denotes the total number of steps in output $i$. These step-wise rewards are also normalized by subtracting their mean and dividing by their standard deviation: $\widetilde{r}_i^{index(j)} = \frac{r_i^{index(j)} - {\rm mean}(\mathbf{R})}{{\rm std}(\mathbf{R})}$. Next, the advantage for each token is computed as the cumulative sum of normalized rewards from all subsequent steps, expressed as $\hat{A}_{i, t} = \sum_{index(j) \ge t} \widetilde{r

\subsubsection{Iterative RL with GRPO} As reinforcement learning training advances, the previously used reward model may no longer adequately guide the current policy model. Hence, we also investigate iterative RL combined with GRPO. As detailed in Algorithm \ref{alg:iter-grpo}, iterative GRPO involves creating new training datasets for the reward model by sampling outputs from the policy model. The existing reward model is then continually updated using a replay mechanism that includes 10\% of past data. Subsequently, we assign the policy model as the reference model and persistently train it with the updated reward model.

\subsection{Training and Evaluating DeepSeekMath-RL}

We perform RL based on DeepSeekMath-Instruct 7B. The RL training data consists of chain-of-thought formatted questions from GSM8K and MATH, drawn from the SFT dataset, totaling approximately 144K questions. Other SFT questions are excluded to specifically assess the effects of RL on benchmarks that have limited data during the RL stage. The reward model training set is constructed following the approach of \citep{wang2023math}. Our initial reward model is trained on DeepSeekMath-Base 7B with a learning rate of 2e-5. For GRPO, the policy model’s learning rate is set at 1e-6, and the KL divergence coefficient is fixed at 0.04. Each question produces $64$ sampled outputs. The maximum sequence length is 1024, and the training batch size is also 1024. The policy model is updated only once after each exploration phase. We assess DeepSeekMath-RL 7B on the same benchmarks as DeepSeekMath-Instruct 7B. For DeepSeekMath-RL 7B, tasks involving GSM8K and MATH with chain-of-thought reasoning are treated as in-domain, while all other benchmarks are considered out-of-domain.

Table \ref{tab:sft_rl_math} presents the results of both open- and closed-source models utilizing chain-of-thought and tool-integrated reasoning on English and Chinese benchmarks. Our observations are: 1) DeepSeekMath-RL 7B achieves accuracies of 88.2\% on GSM8K and 51.7\% on MATH when employing chain-of-thought reasoning. This performance exceeds that of all open-source models in the 7B to 70B parameter range, as well as most closed-source models. 2) Importantly, DeepSeekMath-RL 7B is trained solely on chain-of-thought-format instruction tuning data from GSM8K and MATH, beginning with DeepSeekMath-Instruct 7B. Despite the limited scope of its training data, it outperforms DeepSeekMath-Instruct 7B across all evaluated metrics, demonstrating the effectiveness of reinforcement learning.

\section{Discussion}
In this section, we discuss the insights gained from our pre-training and reinforcement learning experiments.

\subsection{Insights from Pre-Training}

We begin by sharing our experiences related to pre-training. Unless stated otherwise, we follow the training protocols described in Section \ref{sec:quality-policy}. It is important to mention that the DeepSeekMath Corpus referenced here refers to an 89B-token dataset derived from the second iteration of our data collection process.

\subsubsection{Code Training Enhances Mathematical Reasoning}

A widely held yet unproven hypothesis is that training on code improves reasoning capabilities. We aim to provide partial evidence supporting this claim, especially in the context of mathematics: training on code enhances models' proficiency in mathematical reasoning both with and without the use of external tools.

To investigate the impact of code training on mathematical reasoning, we conducted experiments using the following two training regimes:

\textbf{Two-Stage Training}

\begin{itemize}[topsep=0pt]
    \item \textbf{Code Training for 400B Tokens $\rightarrow$ Math Training for 150B Tokens}: DeepSeek-LLM 1.3B is first trained on 400B code tokens, followed by 150B math tokens;
    \item \textbf{General Training for 400B Tokens $\rightarrow$ Math Training for 150B Tokens}: As a comparison, we also conduct experiments using general tokens (drawn from a large-scale general corpus assembled by DeepSeek-AI) instead of code tokens during the initial training phase, aiming to explore the benefits of code tokens over general tokens in enhancing mathematical reasoning.
\end{itemize}

\textbf{One-Stage Training}

\begin{itemize}[topsep=0pt]
    \item \textbf{Math Training for 150B Tokens}: DeepSeek-LLM 1.3B is trained solely on 150B math tokens;
    \item \textbf{Training on a mixture of 400B Code Tokens and 150B Math Tokens}: Since math training after code training reduces coding performance, we examine whether combining code and math tokens in a single-stage training can still improve mathematical reasoning while alleviating the issue of catastrophic forgetting.
\end{itemize}

\paragraph{Results}
Tables \ref{tab:code-math} and \ref{tab:code-math-general-eval} present downstream results under various training configurations.

Code training enhances program-assisted mathematical reasoning in both two-stage and one-stage training paradigms. As depicted in Table \ref{tab:code-math}, within the two-stage training, code training alone significantly boosts the capability to solve GSM8K and MATH problems via Python. Subsequent math training in the second phase brings further gains. Notably, in the one-stage training scenario, integrating code and math tokens effectively addresses the catastrophic forgetting issue observed in two-stage training and simultaneously benefits both coding performance (Table \ref{tab:code-math-general-eval}) and program-assisted mathematical reasoning (Table \ref{tab:code-math}).

Code training also improves mathematical reasoning without relying on tools. In the two-stage approach, the initial code training phase already produces moderate improvements. It additionally enhances the efficiency of later math training, ultimately yielding the best overall performance. However, when code and math tokens are combined in one-stage training, mathematical reasoning without tool use deteriorates. One possible explanation is that DeepSeek-LLM 1.3B’s limited size restricts its ability to fully integrate both coding and mathematical data simultaneously.

\subsubsection{ArXiv Papers Appear Ineffective at Enhancing Mathematical Reasoning}

ArXiv papers are frequently incorporated as part of math pre-training datasets \citep{minerva,gpt-f,llemma,mathpile}. Nevertheless, comprehensive evaluations of their effect on mathematical reasoning remain scarce. Interestingly, contrary to expectations, our experiments indicate that arXiv papers do not effectively improve mathematical reasoning abilities. We test models of varying sizes, including DeepSeek-LLM 1.3B and DeepSeek-Coder-Base-v1.5 7B \citep{deepseek-coder}, trained on arXiv corpora processed through different pipelines:

\begin{itemize}[topsep=0pt]
    \item \textbf{MathPile} \citep{mathpile}: an 8.9B-token corpus curated with cleaning and filtering heuristics, with over 85\% of the content derived from scientific arXiv papers;
    \item \textbf{ArXiv-RedPajama} \citep{redpajama}: a 28.0B-token dataset consisting of the full set of arXiv LaTeX files after removing preambles, comments, macros, and bibliographies.
\end{itemize}

In our study, we independently train DeepSeek-LLM 1.3B for 150B tokens and DeepSeek-Coder-Base-v1.5 7B for 40B tokens on each arXiv dataset. The findings suggest that arXiv papers do not aid in improving mathematical reasoning. Models trained exclusively on arXiv data fail to show significant gains and occasionally exhibit performance declines across various mathematical benchmarks with differing difficulty levels employed here. These benchmarks include quantitative reasoning sets such as GSM8K and MATH (see Table \ref{tab:arxiv-cot}), multiple-choice assessments like MMLU-STEM (Table \ref{tab:arxiv-cot}), and formal mathematics benchmarks like miniF2F (Table \ref{tab:arxiv_atp}).

However, this conclusion has certain caveats and should be interpreted cautiously. The aspects not yet investigated include:

\begin{itemize}[topsep=0pt]
    \item The influence of arXiv tokens on specific mathematics-related tasks excluded from this work, such as the informalization of theorems, i.e., translating formal statements or proofs into informal versions;
    \item The effect of arXiv tokens when used alongside other data types;
    \item Whether any advantages of arXiv papers might emerge at larger model scales.
\end{itemize}

Therefore, additional research is necessary, which we leave to future work.

\subsection{Insights into Reinforcement Learning}
\subsubsection{Towards a Unified Paradigm} This section introduces a unified framework for analyzing diverse training approaches, including SFT, RFT, DPO, PPO, GRPO, and performs experiments to investigate elements within this framework. In general, the gradient with respect to the parameter $\theta$ of a given training method can be expressed as:

\begin{equation}
    \nabla_{\theta}\mathcal{J}_{\textcolor{red}{\mathcal{A}}}(\theta) = \mathbb{E}[\underbrace{(q,o) \sim \textcolor{red}{\mathcal{D}}}_{\text{Source of Data}}]\left( \frac{1}{|o|} \sum_{t=1}^{|o|}  \underbrace{GC_{{\mathcal{A}}}(q, o, t, \textcolor{red}{\pi_{{rf}}})}_{\text{Gradient Coefficient}}  \nabla_{\theta}\log \pi_{\theta}(o_t | q, o_{<t})\right).
\label{eq:objective}
\end{equation}

There are three fundamental elements: 1) \textit{Data Source $\mathcal{D}$}, which provides the data used for training; 2) \textit{Reward Function $\pi_{{rf}}$}, which supplies the training reward signal; 3) \textit{Algorithm $\mathcal{A}$}, which uses the training data and reward signal to generate the gradient coefficient $GC$, controlling the strength of penalty or reward applied to the data. We examine several representative approaches under this common framework:

\begin{itemize}
    \item  \textbf{Supervised Fine-tuning (SFT)}: SFT fine-tunes a pretrained model using human-curated SFT datasets.
    \item \textbf{Rejection Sampling Fine-tuning (RFT)}: RFT further fine-tunes the SFT model on filtered outputs sampled from the SFT model based on SFT questions, selecting outputs according to answer correctness.
    \item \textbf{Direct Preference Optimization (DPO)}: DPO refines the SFT model by fine-tuning it on augmented outputs sampled from the SFT model, utilizing a pairwise DPO loss.
    \item \textbf{Online Rejection Sampling Fine-tuning (Online RFT)}: Unlike RFT, Online RFT starts from the SFT model and iteratively refines it by fine-tuning on augmented outputs sampled from the current policy model in real time.
    \item \textbf{PPO/GRPO}: PPO/GRPO initializes the policy model from the SFT model and enhances it by reinforcing outputs sampled from the live policy model.
\end{itemize}

We present an overview of the components of these approaches in Table \ref{tab:data_policy}. For a more comprehensive derivation, please consult Appendix \ref{app:analysis-rl}.

\paragraph{Observation about Data Source} We categorize the data sources into two types: online sampling and offline sampling. Online sampling refers to training data obtained from the exploration outcomes of the policy model being trained in real-time, whereas offline sampling refers to data derived from the sampling results of the initial SFT model. RFT and DPO adopt the offline sampling approach, while Online RFT and GRPO utilize the online sampling method.

As illustrated in Figure \ref{fig:rl-analysis}, we observe that Online RFT substantially surpasses RFT across two benchmark tests. Specifically, Online RFT performs similarly to RFT in the initial phases of training but gains a clear advantage in later stages, demonstrating the benefits of online training. This observation is intuitive because, early on, the actor and the SFT model are closely aligned, resulting in sampled data with minimal differences. However, as training progresses, data sampled from the actor diverges more significantly, making real-time data sampling more advantageous.

\paragraph{Observation about Gradient Coefficient} The algorithm converts the reward signal into a gradient coefficient to update the model parameters. In our experiments, we classify the reward function as either `Rule' or `Model.' The `Rule' approach assesses response quality based on answer correctness, whereas the `Model' approach involves training a reward model that assigns scores to responses. The reward model’s training data is derived from the rule-based judgments. Equations \ref{eq:GC-RFT} and \ref{eq:GC-GRPO} emphasize a fundamental distinction between GRPO and Online RFT: GRPO specifically modifies its gradient coefficient according to the reward value given by the reward model. This enables the model to differentially reinforce or penalize responses depending on the magnitude of the reward. Conversely, Online RFT does not include this mechanism; it does not penalize incorrect responses and applies uniform reinforcement to all responses that have correct answers, treating them with equal strength.

As illustrated in Figure \ref{fig:rl-analysis}, GRPO outperforms online RFT, underscoring the effectiveness of modifying positive and negative gradient coefficients. Moreover, GRPO+PS demonstrates better results than GRPO+OS, suggesting the advantages of employing fine-grained, step-aware gradient coefficients. Additionally, we investigate iterative RL by performing two rounds of iteration in our experiments. As depicted in Figure \ref{fig:iter-rl}, iterative RL notably enhances performance, particularly during the initial iteration.

\subsubsection{Why Does RL Work?} In this study, we apply reinforcement learning on a subset of instruction tuning data, achieving considerable improvements over the instruction-tuned model. To further clarify why reinforcement learning is effective, we assess the Pass@K and Maj@K accuracy for both the Instruct and RL models across two benchmarks. As shown in Figure \ref{fig:maj-pass}, RL improves Maj@K performance but not Pass@K. These results suggest that RL boosts the model’s overall capability by making the output distribution more robust; in other words, \textbf{the enhancement appears to stem from increasing the probability of the correct answer appearing within the TopK outputs rather than improving fundamental capabilities.} Correspondingly, \citep{wang2023making} identified a \textbf{misalignment issue} in reasoning tasks within the SFT model and showed that reasoning performance can be enhanced by employing a series of preference alignment techniques \citep{yuan2023rrhf,song2023preference,wang2023making}.

\subsubsection{How Can More Effective RL Be Achieved?}
\label{sec:effec_rl}
We demonstrate that RL is quite effective in mathematical reasoning tasks and present a unified framework for understanding various representative training methods. Within this framework, all approaches can be viewed as either direct or simplified RL techniques. As outlined in Equation \ref{eq:objective}, there are three main components: Data Source, Algorithm, and Reward Function. We suggest several promising future directions relating to these components.

\paragraph{Data Source} The data source serves as the foundational input for all training methods. Specifically, in RL, this refers to unlabeled questions with outputs sampled from the policy model. In this work, we utilize only the questions from the instruction tuning stage and apply a basic nucleus sampling method to generate outputs. We believe this might partly explain why our RL pipeline improves only Maj@K performance. Future work will explore the RL pipeline on out-of-distribution question prompts combined with \textbf{advanced sampling (decoding) strategies}, such as those based on tree-search methods \citep{yao2023tree}. Additionally, \textbf{efficient inference techniques} \citep{xia-etal-2023-speculative,leviathan2023fast,kwon2023efficient,xia2024unlocking}, which influence the exploration efficiency of policy models, will also be crucial.

\paragraph{Algorithms} Algorithms utilize the data and reward signals to compute the gradient coefficients, which are then used to update the model parameters. According to Equation \ref{eq:objective}, all current methods to some degree fully \textbf{TRUST} the reward function's signal to either increase or decrease the conditional probability of a specific token. Nonetheless, it is unrealistic to guarantee that the reward signal is consistently reliable, particularly in highly complex tasks. For instance, the PRM800K datasets \citep{lightman2023let}, despite being meticulously annotated by skilled annotators, still contain roughly 20\% incorrect annotations\footnote{\url{https://github.com/openai/prm800k/issues/12\#issuecomment-1728491852}}. Therefore, we aim to investigate reinforcement learning algorithms that are robust to noisy reward signals. We posit that such \textbf{WEAK-TO-STRONG} \citep{burns2023weak} alignment methods will fundamentally transform learning algorithms.

\paragraph{Reward Function} The reward function serves as the origin of the training signal. In reinforcement learning, the reward function is typically implemented as a neural reward model. We identify three crucial directions for the development of reward models: 1) \textbf{Enhancing the generalization capability of the reward model.} The reward model must generalize effectively to out-of-distribution questions and sophisticated decoding outputs; otherwise, reinforcement learning might only stabilize the LLM distribution without improving its core capabilities; 2) \textbf{Capturing the uncertainty of the reward model.} This uncertainty could serve as a bridge linking weak reward models and weak-to-strong learning algorithms; 3) \textbf{Efficiently constructing high-quality process reward models} that can deliver fine-grained training signals for the reasoning process \citep{lightman2023let,wang2023math}.

\section{Conclusion, Limitation, and Future Work} We introduce DeepSeekMath, which surpasses all open-source models on the competition-level MATH benchmark and nearly matches the performance of closed models. DeepSeekMath~is initialized from DeepSeek-Coder-v1.5 7B and undergoes continuous training over 500B tokens, with a large portion of the training data comprising 120B math tokens extracted from Common Crawl. Our comprehensive ablation study reveals that web pages provide considerable promise for obtaining high-quality mathematical data, whereas arXiv appears less beneficial than initially anticipated. We propose Group Relative Policy Optimization (GRPO), a variant of Proximal Policy Optimization (PPO), which significantly enhances mathematical reasoning abilities while using less memory. Experimental results demonstrate that GRPO remains effective even after DeepSeekMath-Instruct 7B attains high benchmark scores. Additionally, we offer a unified framework to interpret a series of methods and highlight several promising avenues for more efficient reinforcement learning.

While DeepSeekMath~achieves remarkable results on quantitative reasoning benchmarks, its performance in geometry and theorem-proof tasks is comparatively weaker than that of closed models. For example, in preliminary tests, the model struggles with problems involving triangles and ellipses, suggesting possible data selection biases during pre-training and fine-tuning. Furthermore, limited by model scale, DeepSeekMath~underperforms GPT-4 in few-shot learning. GPT-4 can enhance its performance with few-shot examples, whereas DeepSeekMath~exhibits similar outcomes in both zero-shot and few-shot evaluations. Moving forward, we plan to refine our data selection pipeline to create a more curated and high-quality pre-training corpus. Additionally, we intend to investigate the potential directions outlined in Section \ref{sec:effec_rl} to achieve more effective reinforcement learning for LLMs.

\end{document}